\newcommand{\verfasser}{Phil Steinhorst}					% Verfasser des Dokuments
\newcommand{\fach}{Elliptische Kurven und Kryptographie}				% Titel der Vorlesung								
\newcommand{\shortFach}{EKK}								% Titel kurz
\newcommand{\prof}{PD Dr. Karin Halupczok}					% Dozent
\newcommand{\untertitel}{gelesen von \prof}					% Untertitel
\newcommand{\semester}{Sommersemester 2015}					% Semester
\newcommand{\homepage}{http://wwwmath.uni-muenster.de/u/karin.halupczok/ellKKSoSe15/}	% Vorlesungshomepage
% % % % % % % % % % % % % % % % % % % % % % % % % % % % % % % % %
%		PhistScript.tex											%
%		XeLaTeX-Settings for Scripts							%
%																%
%		Author: Phil Steinhorst									%
% % % % % % % % % % % % % % % % % % % % % % % % % % % % % % % % %

\setlength\parindent{0pt}             % ohne Einrueckung

% % % Konfiguration von scrheadings
\setheadsepline{1pt}[\color{black}]
\pagestyle{scrheadings}
\clearscrheadfoot

% % % Kopf-/Fußzeilenlayout
\providecommand{\shortFach}{\fach}
\ohead{\small \leftmark}
\ofoot[{ \small \thepage}]{ \small \thepage} 
\automark{section}

% Metadaten
\title{\fach}
\author{\verfasser}
\subtitle{\untertitel}
\date{\semester}

% Inhaltsverzeichnis
\setcounter{tocdepth}{1}
\usepackage[tocindentauto]{tocstyle}
\usetocstyle{KOMAlike}	

% Nummerierung
\numberwithin{equation}{section}		% Section in Equation-Nummerierung einbeziehen

% ntheorem package
\newcounter{tcount}						% Zähler für Theoreme
\numberwithin{tcount}{section}			% Section in Theorem-Nummerierung einbeziehen
\usepackage[hyperref]{ntheorem}			% Lade ntheorem mit hyperref-Workaround
\theoremstyle{break}					% Zeilenumbruch nach Theoremkopf
\theorembodyfont{\normalfont}			% nicht kursiv
\theorempreskipamount0.5cm				% Abstand vor Theorem
\theorempostskipamount0.5cm				% Abstand nach Theorem
\newtheorem{thm}[tcount]{Theorem}
\newtheorem{bsp}[tcount]{Beispiel}
\newtheorem{defn}[tcount]{Definition}
\newtheorem{bem}[tcount]{Bemerkung}
\newtheorem{lemma}[tcount]{Lemma}

% Workaround für Tabulatoren im Theorem-Verzeichnis
\makeatletter
\def\thm@@thmline@name#1#2#3#4{%
        \@dottedtocline{-2}{0em}{2.3em}%
                   {\makebox[\widesttheorem][l]{#1 \protect\numberline{#2}}#3}%
                   {#4}}
\@ifpackageloaded{hyperref}{
\def\thm@@thmline@name#1#2#3#4#5{%
    \ifx\\#5\\%
        \@dottedtocline{-2}{0em}{2.3em}%
            {\makebox[\widesttheorem][l]{#1 \protect\numberline{#2}}#3}%
            {#4}
    \else
        \ifHy@linktocpage\relax\relax
            \@dottedtocline{-2}{0em}{2.3em}%
                {\makebox[\widesttheorem][l]{#1 \protect\numberline{#2}}#3}%
                {\hyper@linkstart{link}{#5}{#4}\hyper@linkend}%
        \else
            \@dottedtocline{-2}{0em}{2.3em}%
                {\hyper@linkstart{link}{#5}%
                  {\makebox[\widesttheorem][l]{#1 \protect\numberline{#2}}#3}\hyper@linkend}%
                    {#4}%
        \fi
    \fi}
}
\makeatother
\newlength\widesttheorem
\AtBeginDocument{
  \settowidth{\widesttheorem}{Proposition 10.10\quad}	% Referenzstring für die Breite der ersten Spalte
}

\theoremlisttype{optname}	% nur Theoreme auflisten, die einen Namen haben. 							% XeLaTeX-Konfiguration für Skripte
\usepackage{easybmat}
% % % Dokumentspezifische Einstellungen % % %
\setcounter{tocdepth}{3}				% Tiefe im Inhaltsverzeichnis (1: nur Sections, 2: auch subsections...)

\numberwithin{equation}{section}			% Section in Equation-Nummerierung einbeziehen
\newcounter{counter}						% Zähler für Sätze
\newcounter{lecture}
\numberwithin{counter}{lecture}		% Section in Satz-Nummerierung einbeziehen

\newcommand{\nextlecture}{\stepcounter{lecture} \setcounter{counter}{0}}


% ab hier können neue Umgebungen für Sätze definiert und benannt werden!
% Syntax: \newtheorem{interner Umgebungsname}[Zähler]{Name im Dokument}
% Möchte man Sätze und Definitionen getrennt durchnummerieren, benötigt man weitere Zähler.
\theoremlisttype{optname}					% nur Theoreme auflisten, die einen Namen haben. 
\newtheorem{satz}[counter]{Satz}
\newtheorem{defn}[counter]{Definition}
\newtheorem{bem}[counter]{Bemerkung}
\newtheorem{bsp}[counter]{Beispiel}
\newtheorem{lemma}[counter]{Lemma}
\newtheorem{folg}[counter]{Folgerung}
\newtheorem{zusatz}[counter]{Zusatz}
\newtheorem{anw}[counter]{Anwendung}
\newtheorem{erl}[counter]{Erläuterung}
\newtheorem{mot}[counter]{Motivation}
\newtheorem{bew}[counter]{Beweis}
\newtheorem{kor}[counter]{Korollar}
\newtheorem{vor}[counter]{Vorbereitung}
% % % % % % % % % % % % % % % % % % % % % % %
\newcommand{\PP}{\mathbb{P}}
\DeclareMathOperator{\li}{li}
\newcommand{\kon}{\equiv}
\newcommand{\oh}{\mathcal{O}}
\newcommand{\cc}{\mathcal{C}}
\DeclareMathOperator{\Res}{Res}
\DeclareMathOperator{\disc}{disc}
\DeclareMathOperator{\Rg}{Rg}
\newcommand{\pol}{\mathcal{P}}
\newcommand{\leg}[2]{\enbrace*{\frac{#1}{#2}}}

% % % % % % % % % % % % % % % % % % % % % % %
\begin{document}
\begin{titlepage}
	% Nach einer Vorlage von http://www.LaTeXTemplates.com
	\newcommand{\HRule}{\rule{\linewidth}{0.8mm}} % Defines a new command for the horizontal lines, change thickness here

	\center % Center everything on the page
 
	\begin{minipage}{0.4\textwidth}
	\begin{flushleft}
	\includegraphics[height=1.5cm,keepaspectratio]{img/wwulogo.pdf}\\[1cm]
	\end{flushleft}
	\end{minipage}
	\hfill
	\begin{minipage}{0.4\textwidth}
	\begin{flushright}
	\vspace*{0.3cm}
	\includegraphics[height=1.2cm,keepaspectratio]{img/fb10logo.pdf} \
	\end{flushright}
	\end{minipage}

	\vspace{2cm}
	


	\HRule \\[0.8cm]
	{ \huge \sffamily\bfseries $\mathbf{CAT}\boldsymbol{(0)}$ kubische Komplexe}\\[0.4cm] % Title of your document
	\HRule \\[1cm]
 
	{\LARGE gelesen von} \\[.7cm]
	\textsc{\LARGE \textbf{Dr. Olga Varghese}}\\[.7cm]
	{\LARGE im Wintersemester 2015/2016}\\[2cm]

	\begin{tikzpicture}[scale=2]
		\draw[very thick,dashed] (1,1) -- (1.45,1.45) -- (2.45,1.45) -- (2,1);
		\draw[very thick,dashed] (2.45,1.45) -- (2.45,1);
		\draw[very thick,dashed] (1.45,1.45) -- (1.45,.45) -- (1,0);
		\draw[very thick,dashed] (1.45,.45) -- (2.45,.45);
		\draw[very thick,dashed] (3,1) -- (3.45,1.45) -- (4.45,1.45);
		\draw[very thick,dashed] (3.45,1.45) -- (3.45,2.45);
				
		\draw[very thick] (0,1) -- (2,1) -- (2,0) -- (2.45,.45) -- (2.45,1);
		\draw[very thick] (2,0) -- (1,0) -- (1,2) -- (4,2) -- (4,1) -- (3,1) -- (3,3) -- (2,3) -- (2,0);
		\draw[very thick] (2,1) -- (3,1) -- (3,0) -- (5,0);
		\draw[very thick] (4,0) -- (4,1) -- (4.45,1.45) -- (5.45,1.45) -- (5.45,2.45) -- (3.45,2.45) -- (3,2);
		\draw[very thick] (4,2) -- (4.45,2.45) -- (4.45,1.45);
		\draw[very thick] (4.45,2.45) -- (4.45,3.45);

		\draw[ultra thick,color=red,dashed] (3,1.5) -- (3.45,1.95) -- (4.45,1.95);		
		\draw[ultra thick,color=red] (1,1.5) -- (4,1.5) -- (4.45,1.95) -- (5.45,1.95);
	\end{tikzpicture}

	\vfill
	{\Large Vorlesungsmitschrift von Phil Steinhorst} \\[.5cm]
	{\large Stand: \today}
	
\end{titlepage}							% Titelseite ausgeben
\section*{Vorwort}
\label{sec:preface}
	Der vorliegende Text ist eine inhaltliche Aufbereitung zur Vorlesung Elliptische Kurven und Kryptographie, gelesen von PD Dr. Karin Halupczok an der WWU Münster im Sommersemester 2015. Der Inhalt entspricht weitestgehend dem handschriftlichen Skript, welches auf der Vorlesungswebsite bereitsgestellt wird. Dieses Werk ist daher keine Eigenleistung des Autors und wird nicht von der Dozentin der Veranstaltung korrekturgelesen. Für die Korrektheit des Inhalts wird keinerlei Garantie übernommen. Bemerkungen, Korrekturen und Ergänzungen kann man folgenderweise loswerden:
	\begin{itemize}
		\item persönlich durch Überreichen von Notizen oder per E-Mail
		\item durch Abändern der entsprechenden TeX-Dateien und Versand per E-Mail an mich
		\item direktes Mitarbeiten via GitHub. Dieses Skript befindet sich im \texttt{latex-wwu}-Repository von Jannes Bantje:
		\begin{center}
			\url{https://github.com/JaMeZ-B/latex-wwu}
		\end{center}
	\end{itemize}

\subsection*{Literatur}
\label{sub:lit}
\begin{itemize}
	\item Blake, Seroussi, Smart: Elliptic curves in cryptography
	\item Menezes, van Oorschot, Vanstone: Handbook of applied cryptography
	\item Silverman: The arithmetic of elliptic curves
	\item Silverman: A friendly introduction to number theory, chap. 40-45
	\item Washington: Elliptic curves, number theory and cryptography
	\item Werner: Elliptische Kurven in der Kryptographie
\end{itemize}

\subsection*{Kommentar der Dozentin}
In der Vorlesung beschäftigen wir uns mit den arithmetischen und geometrischen Eigenschaften elliptischer Kurven sowie deren Anwendungen in der Kryptographie. Dabei werden wir auch einen Vergleich mit Anwendungen der elementaren Zahlentheorie in der Kryptographie ziehen. Wir verfolgen eine elementare Herangehensweise, d.h. Kenntnisse der algebraischen Geometrie und der Funktionen- oder Zahlentheorie werden nicht benötigt. Es genügen die Vorkenntnisse aus den Grundvorlesungen.

\subsection*{Vorlesungswebsite}
\label{sub:link}
Das handgeschriebene Skript sowie weiteres Material findet man unter folgendem Link:
\begin{center}
	\url{\homepage}
\end{center}

\subsection*{Titelbild}
\label{sub:titlepic}
Das fehlt noch. Über Ideen und Anregungen freue ich mich sehr!

\vfill
\begin{flushright}
	Phil Steinhorst \\
	p.st@wwu.de
\end{flushright}
\newpage
\tableofcontents
\newpage											% Inhaltsverzeichnis ausgeben
\section{Motivation und Einführung}
\label{sec:para0}

Motivation und Einführung. \marginnote{13.04. \\ \ [1]}
\newpage
\chapter{CAT(0)-Räume} % (fold)
\label{cha:1}

\begin{definition}[metrischer Raum]
	Sei $X$ eine Menge. Eine Abbildung $d \colon X \times X \rightarrow \RR_{\geq 0}$ heißt \Index{Metrik}, wenn für alle $x,y,z \in X$ gilt:
	\begin{enumerate}[(i)]
		\item $d(x,y) = 0 \Leftrightarrow x=y$
		\item $d(x,y) = d(y,x)$
		\item $d(x,z) \leq d(x,y) + d(y,z)$
	\end{enumerate}
	Das Paar $(X,d)$ heißt dann \Index{metrischer Raum}.
\end{definition}

\cleardoubleoddemptypage
\section{Der euklidische Algorithmus}
\label{sec:para2}
	Sei $R$ kommutativer Ring mit $1 \neq 0$. Für beliebiges $a \in R$ betrachte man die Menge der Vielfachen von $a \in R$, also
	\[ Ra := {xa : x \in R} = {b \in R : a|b} \]
	Die Teilmenge $I = Ra$ hat folgende Eigenschaften:
	\begin{enumerate}[(i)]
		\item $0 \in I$
		\item $b_1,b_2 \in I \Rightarrow b_1+b_2 \in I$
		\item $c \in R, b \in I \Rightarrow cb \in I$
	\end{enumerate}
	
\begin{defn}[Ideal, Hauptideal] \label{def_2.1}
	Eine Teilmenge $I$ von $R$ heißt ein \Index{Ideal} in $R$, falls die Eigenschaften (i), (ii), (iii) erfüllt sind. $I$ heißt \Index{Hauptideal}, wenn es ein $a \in R$ gibt mit $I = Ra$. Wir verwenden die Bezeichnung
	\[ (a) := Ra \]
	und nennen $(a)$ das von $a \in R$ erzeugte Hauptideal.
\end{defn}

\minisec{Bemerkung}
	\begin{enumerate}[(1)]
		\item $(b) \subseteq (a) \Leftrightarrow a | b$
		\item $a \assoz b \Leftrightarrow (a) = (b)$
		\item $c$ ist gemeinsames Vielfaches von $a_1,\dots,a_n \Leftrightarrow (c) \subseteq (a_1) \cap \dots \cap (a_n)$
		\item $m$ ist ein kgV von $a_1,\dots,a_n \Leftrightarrow (a_1) \cap \dots \cap (a_n) = (m)$
		\item $d$ ist ein gemeinsamer Teiler von $a_1,\dots,a_n \Leftrightarrow (a_i) \subseteq (d)$ für $1 \leq i \leq n$
		\item $d$ ist ein gemeinsamer Teiler von $a_1,\dots,a_n \Leftrightarrow Ra_1 + Ra_2 + \dots + Ra_n \subseteq (d)$
		\item $d$ ist ein ggT von $a_1,\dots,a_n \Leftrightarrow (d)$ ist das kleinste Hauptideal mit $Ra_1 + \dots Ra_n \subseteq (d)$.
	\end{enumerate}
	Ein ggT lässt sich also idealtheoretisch nicht so einfach charakterisieren wie oben ein kgV durch (4). Am schönsten wäre es, wenn $Ra_1 + \dots + Ra_n$ ein Hauptideal wäre, dann würde (7) übergehen in:
	\[ d \text{ ist ein ggT von } a_1,\dots,a_n \Leftrightarrow Ra_1+Ra_2+ \dots + Ra_n = (d) \]
	
\begin{defn}[Hauptidealring] \label{def_2.2}
	Ein Integritätsring $R$ heißt ein \Index{Hauptidealring}, wenn jedes Ideal $I$ von $R$ ein Hauptideal ist.
\end{defn}

\minisec{Bezeichnung}
	Für Elemente $a_1,\dots,a_n$ in einem beliebigen kommutativen Ring $R$ mit $1 \neq 0$ setze
	\[ (a_1,\dots,a_n) := Ra_1 + \dots + Ra_n \]
	Man nennt $(a_1,\dots,a_n)$ das von $a_1,\dots,a_n$ erzeugte Ideal in $R$.
\newpage
\section{Partielle Differentialgleichungen erster Ordnung}
\label{sec:para3}

% % % % % % 22. Apr.
\minisec{Motivation: Linearer Fall}
	Betrachte \marginnote{22. Apr}
	\[ a(x) \cdot \nabla u(x) + a_0(x) = 0 \] mit Randdaten auf glatter $(n-1)$-dimensionaler Hyperfläche $\Gamma \subset \RR^n$, $u(x) = g(x)$ auf $\Gamma$. \\
	Die Richtungsableitung von $u$ in Richtung $a$ ist $a(x) \cdot \nabla u(x) = -a_0 (x)$, d.h. wir wissen, wie sich $u$ entlang charakteristischer Kurven $s \mapsto x(s)$ ändert mit
	\begin{equation}
	\begin{aligned}
		\dot{x}(s) = a(x(s)) \label{eq_2204_1}
	\end{aligned}
	\end{equation} \\
	Sei $\Gamma$ parametrisiert durch $x_0 \colon \RR^{n-1} \rightarrow \Gamma$. Zu $\widehat{x} \in \RR^n$ können wir $s \in \RR, y \in \RR^{n-1}$ finden, sodass $\widehat{x} = x(s)$ für eine Lösung $x(s)$ von \eqref{eq_2204_1} mit $x(0) = x_0(y)$, d.h. wir reparametrisieren $\RR^n$ durch eine Funktion $\widehat{x}\colon \RR^{n-1} \times \RR \rightarrow \RR, (y,s) \mapsto \widehat{x}(y,s)$. Dann gilt:
	\[ u(\widehat{x}(y,s)) = \underbrace{u(\widehat{x}(y,0))}_{g(\widehat{x}(y,0))} + \int\limits_{0}^{s} \underbrace{\frac{\der u(\widehat{x}(y,t))}{\der \tau}}_{=-a_0(\widehat{x}(y,\tau))} d\tau; \quad u(x) = (u \circ \widehat{x}) \circ \widehat{x}^{-1} \]
	
\begin{bem} \label{bem_2}
	Damit $u$ auf $\RR^n$ definiert ist, muss $\widehat{x}^{-1}$ definiert sein.\marginnote{[2]} Für lokale Invertierbarkeit ist hinreichend:
	\[ \det(D \widehat{x}) \neq 0 \]
	Auf $\Gamma$ ist $\det(D \widehat{x}(y,0) = \det( \underbrace{D x_0(y)}_{\in R^{n\times m -1}}\vert a(x_0(y)))$, d.h. wir brauchen $a(x_0(y)) \notin \operatorname{span}(Dx_0(y))$. Die Spalten von $Dx_0(y)$ spannen den Tangentialraum von $\Gamma$ an $x_0(y)$ auf. Damit dies ungleich 0 ist, muss $a \notin \operatorname{span}(Dx_0)$, d.h. $a \cdot \nu \neq 0$ mit $\nu$ als Normale auf $\Gamma$, d.h. $\Gamma$ ist nichtcharakteristisch.
\end{bem}
	
\begin{bsp} \label{bsp_3}
	Betrachte \marginnote{[3]}
	\[ \begin{cases}
		u_{x_1} + u_{x_2} = 1 \text{ auf } \RR^2 \\
		u = 0 \text{ auf } \Gamma = \{x_1 + x_2 = 0 \}
	\end{cases} \]
	Charakteristische Kurven: $\dot{x} = \begin{pmatrix} 1 \\ 1	\end{pmatrix} \Rightarrow x(s) = \begin{pmatrix} c_1 + s \\ c_2 + s \end{pmatrix}$. \\
	Parametrisiere $\Gamma$:
	\[ x_0(y) := \begin{pmatrix} y \\ -y \end{pmatrix} = x_0(y) \]
	$\Rightarrow \widehat{x}(y,s) = x(s)$ für die Kurve $x(s)$ mit $x(0) = x_0(y) \Rightarrow \widehat{x}(y,s) = \begin{pmatrix} y+s \\ -y+s \end{pmatrix}$ \\
	$\Rightarrow \widehat{x}^{-1}(x) = (\frac{x_1-x_2}{2},\frac{x_1+x_2}{2})$ \\
	$\Rightarrow u(\widehat{x}(y,s)) = 0 + \int\limits_{0}^{s} 1 d\tau = s$ \\
	$\Rightarrow u(x) = \underbrace{(u \circ \widehat{x})}_s \circ \widehat{x}^{-1}(x) = \frac{x_1+x_2}{2}$
\end{bsp}
	
\minisec{Nichtlinearer Fall}
	Betrachte
	\begin{equation}
		0 = F(\nabla u(x),u(x),x), \quad u=g \text{ auf } \Gamma, \quad F,g,\Gamma \ C^2\text{-glatt} \label{eq_2204_2}
	\end{equation}
	Idee: Information fließt entlang charakteristischer Kurven \\
	Ziel: Finde charakteristische Kurven und wandele die PDGL um in ein System von gewöhnlichen DGL entlang der Kurven. \\
	Notation: \begin{itemize}
	\item Für eine Kurve $x(s)$ schreibe $z(s) := u(x(s))$ und $p(s) = \nabla u(x(s))$, d.h.
	\[ 0 = F(p(s), z(s), x(s)) \]
	\item $\dot{} = \frac{\der}{\der s}$
	\item $F_p,F_z,F_x$ sind partielle Ableitungen von $F$ nach erstem, zweitem und drittem Argument.
	\end{itemize}
	Suche GDGL für $x,p,z$.
	
\begin{thm}[Charpits Gleichungen]
\label{thm_charpits} \label{thm_4}
	Sei $u \in C^2(\Omega)$ eine Lösung von \eqref{eq_2204_2}.\marginnote{[4]} Entlang von Kurven $x(s)$ in $\Omega$ mit
	\begin{equation}
	\begin{aligned}
		\dot{x} = F_p \label{eq_2204_3}
	\end{aligned}
	\end{equation}
	gilt
	\begin{align}
		&\dot{z} = p \cdot F_p \label{eq_2204_4} \\
		&\dot{p} = -pF_z - F_x \label{eq_2204_5}
	\end{align}
	Diese GDGLs heißen \Index{Charpits} oder \Index{charakteristische Gleichungen}, Kurven $x(s)$ oder $(x(s),z(s),p(s))$ heißen \Index{charakteristische Kurven}.
\end{thm}
	
\minisec{Beweis}
	Betrachte beliebige Kurve $x(s)$. Dann gilt:
	\[ \begin{array}{c}
		\dot{z}(s) = \nabla u(x(s)) \cdot \dot{x}(s) = p(s) \cdot \dot{x}(s) \\
		\dot{p}(s) = D^2u(x(s)) \dot{x}(s)
	\end{array} \]
	Um $D^2u$ loszuwerden, differenziere \eqref{eq_2204_2} nach $x$:
	\[ 0 = \frac{d}{dx} F(\nabla u(x),u(x),x) = D^2 u F_p + F_z \overbrace{\nabla u}^p + F_x \]
	Wähle die Kurve $x(s)$, sodass $\dot{x}(s) = F_p(p(s),z(s),x(s))$, so erhalten wir
	\[ \dot{p}(s) = -F_z p - F_x \qquad \dot{z} = p \cdot F_p \]. \qed

\mbox{} \\
\eqref{eq_2204_3}, \eqref{eq_2204_4} und \eqref{eq_2204_5} können mit Anfangswerten $x(0),z(0),p(0)$ gelöst werden. Sei $\Gamma$ parametrisiert durch $x_0(y), y \in \RR^{n-1}$. Werte für $x(0)$ und $z(0)$ sind auf $\Gamma$ gegeben durch $x(0) = x_0(y) \in \Gamma$ und $z(0) = g(x_0(y))$. $p(0)$ erhalten wir durch Lösen von
\begin{equation}
	\left. \begin{cases}
	0 = F(p(0),z(0),x(0)) & \text{PDGL} \\
	0 = \nabla x_0 (y) p(0) - \nabla_y g(x_0(y)) & \text{Abl. von } u(x_0(y)) = g(x_0(y)) \text{ nach } y
	\end{cases} \right\}  \label{eq_2204_7} \end{equation}
wobei
	\[ \nabla x_0 (y) = \begin{pmatrix}
	\der_{y_1} {x_0}_1 & \cdots & \der_{y_1} {x_0}_n \\ 
	\vdots & \ddots & \vdots \\ 
	\der_{y_{n-1}} {x_0}_1 & \cdots & \der_{y_{n-1}} {x_0}_n
	\end{pmatrix} \]
Die Lösung von \eqref{eq_2204_3}, \eqref{eq_2204_4} und \eqref{eq_2204_5} für diese Anfangswerte liefert eine Parametrisierung $\widehat{x}(y,s)$ des $\RR^n$ mit $\widehat{x}(y,s) = x(s)$ für die Kurve $x(s)$ mit $x(0) = x_0(y)$ sowie Funktionen $z(y,s),p(y,s)$.

\begin{defn}[nichtcharakteristische Hyperfläche] \label{def_nichtchar} \label{def_5}
	Eine glatte $(n-1)$-dimensionale Hyperfläche\marginnote{[5]} $\Gamma \subset \RR^n$ heißt \Index{nichtcharakteristisch} zu \eqref{eq_2204_2}, wenn für die Einheitsnormale $\nu$ an $\Gamma$ gilt:
	\[ F_p(p,z,x) \cdot \nu(x) \neq 0 \text{ für alle } x \in \Gamma, z \in \RR, p \in \RR^n \]
\end{defn}
	
\begin{lemma} \label{lemma_6}
	Wenn gilt: \marginnote{[6]} \begin{itemize}
	\item $\Gamma$ ist nichtcharakteristisch
	\item Für ein gegebenes $y_0 \in \RR^{n-1}$ existiert eine Lösung von \eqref{eq_2204_7}
	\end{itemize}
	dann: \begin{itemize}
		\item kann \eqref{eq_2204_7} auch in einer Umgebung $U$ von $y_0$ nach $p(0)$ gelöst werden
		\item existieren ein offenes Intervall $I = (-\delta,\delta)$ und eine Umgebung $V$ von $x_0(y_0)$, sodass für alle $x \in V$ ein $y \in U$ und $s \in I$ existiert mit $x = \widehat{x}(y,s)$. $\widehat{x}^{-1}$ ist $C^2$.
	\end{itemize}
\end{lemma}
	
\minisec{Beweis}
	\begin{itemize}
		\item Existenz von $p(0)$ folgt direkt aus dem Satz über implizite Funktionen. (Beachte:
		\[ \frac{\der}{\der p(0)} \begin{pmatrix}
		F(p(0),z(0),x(0)) \\ \nabla x_0 (y) p(0) - \nabla_y g(x_0(y)) \end{pmatrix} = \begin{pmatrix}
		F_p \\ \nabla x_0 (y) \end{pmatrix} \]
		hat vollen Rang.)
		\item Existenz und Glattheit von $\widehat{x}^{-1}$ folgt aus dem Satz über inverse Funktionen, da $\det (y_0,0) = \det \begin{pmatrix}	\nabla x_0 \\ F_p \end{pmatrix} \neq 0$ \qed
	\end{itemize}
	
\begin{bem} \label{bem_8}
	Auch hier bedeutet nichtcharakteristisch,\marginnote{[7]} dass auf $\Gamma$ nach den partiellen Ableitungen von $u$ gelöst werden kann.
\end{bem}
	
\begin{thm}[Lokale Existenz einer glatten Lösung]
\label{thm_lok_existenz} \label{thm_8}
	Seien \marginnote{[8]} \begin{itemize}
		\item $\Gamma, g, F$ glatt, also $C^2$
		\item $\Gamma$ nichtcharakteristisch und
		\item es gebe ein $y_0 \in \RR^{n-1}$ und ein $p(0) \in \RR^n$, das \eqref{eq_2204_7} löst.
	\end{itemize}
	Dann existiert auf der Umgebung $V$ von $x_0(y_0)$ aus Lemma \ref{lemma_6} eine $C^2$-Lösung $u$ von \eqref{eq_2204_2} und \eqref{eq_2204_3}. Es ist $u(x) = z(y,s)$ und $\nabla u(x) = p(y,s)$ für $(y,s) = \widehat{x}^{-1} (x)$.
\end{thm}

% % % % % % % % % % % % % 25.04.
\minisec{Beweis}
	\begin{enumerate}[1.]
		\item Nach Lemma \eqref{lemma_6} \marginnote{25. Apr} existiert $\widehat{x}^{-1}$ auf $V$ sowie $z(y,0) = g(x_0(y)),p(y,0)$ mit $F(p(y,0),z(y,0),\widehat{x}(y,0)) = 0$ und $\nabla x_0 (y) p(0) - \nabla_y g(x_0(y)) = 0$. Löse Charpits Gleichungen mit diesen Anfangswerten $\Rightarrow z(y,s),p(y,s)$.
		\item Zeige: $f(y,s) := F(p(y,s),z(y,s),\widehat{x}(y,s)) = 0$. \\
			$f(y,0) = 0, \dot{f}(y,s) = F_p \cdot \dot{p} + F_z \cdot \dot{z} + F_x \dot{\widehat{x}} = F_p(-pF_z - F_x) + F_zpF_p + F_x \cdot F_p = 0$
		\item Zeige $p(\widehat{x}^{-1}(x)) = \nabla u(x)$ für $u(x) = z(\widehat{x}^{-1}(x))$. \marginnote{$z \equiv u \circ \widehat{x}^{-1}$ \\ $z_s = \nabla u \cdot \widehat{x}_s = p \cdot \widehat{x}_s$ \\ $z_y = \nabla u \widehat{x}_y = p \cdot \widehat{x}_y$} \begin{itemize}
				\item $\cdot{z}(y,s) = p(y,s) \cdot \dot{\widehat{x}}(y,s)$
				\item $z_y(y,s) = p^T(y,s) \widehat{x}_y(y,s)$, denn $r(s) := z_y(y,s) - p^T(y,s) \widehat{x}_y (y,s)$ erfüllt
				\begin{equation}
				\begin{aligned}
					\dot{r}(s) &= \underbrace{\dot{z_y}(y,s)}_{p^T_y \dot{\widehat{x}} + p^T \dot{\widehat{x}}_y} - \dot{p}^T(y,s) \widehat{x}_y(y,s) - p^T(y,s) \dot{\widehat{x}}_y(y,s) = p_y^T F_p - (-pF_z - F_x)^T \widehat{x}_y \\ \notag
					&= \frac{\der}{\der y} \underbrace{F(p(y,s),z(y,s),\widehat{x}(y,s))}_{=0} + \underbrace{p^T F_z \widehat{x}_y - F_z z_y}_{-F_z r}
				\end{aligned}
				\end{equation} 
				mit Anfangswert $r(0) = \underbrace{z_y(y,0)}_{\nabla_y g(x_0(y))} - p^T(y,0) \underbrace{\widehat{x}_y(y,0)}_{\nabla_y x_0(y)} = 0$. $\Rightarrow r \equiv 0$.
				\item $u_x = z_y y_x + z_s s_x = p^T \widehat{x}_y y_x + p^T \widehat{x}_s s_x = p^T \underbrace{(\widehat{x}_y y_x + \widehat{x}_s s_y)}_{\widehat{x}_x} = p^T, p = \nabla u$. \qed
			\end{itemize}
	\end{enumerate}

\begin{bsp} \label{bsp_9}
	Löse das Cauchy-Problem \marginnote{[9]}
	\[ \begin{cases}
		u_t + uu_x = 1 & \text{ auf } (0,\infty) \times \RR \\
		u = x & \text{ auf } \Gamma = \{t = 0\} \end{cases} \]
	Charakteristische Kurven sind gegeben durch $\frac{\der}{\der s} \begin{pmatrix} (t,x) \\ z \\ 	p \end{pmatrix} = \begin{pmatrix} (1,z) \\ P-1+zp_2 \\ -p_2 p \end{pmatrix}$, also
	\[ \begin{pmatrix} (t,x) \\ z \\ 	p \end{pmatrix} = \begin{pmatrix} (t(0)+s,x(0) + z(0)s + \frac{p_1(0)+z(0)p_2(0)}{2} s^2 \\ p_1(0)s + z(0)(1+p_2(0)s) \\ \frac{p}{1+p_2(0)s} \end{pmatrix}.\]
	Parametrisiere $\Gamma$ durch $x_0(y) = (0,y)$. Das liefert Anfangswerte $(t(0),x(0))=x_0(y)=(0,y)$ und $z(0) = y$. Lösen von \eqref{eq_2204_7} liefert $p(0) = \begin{pmatrix} 1-y \\ 1 \end{pmatrix}$. Insgesamt ergibt sich:
	\[ \Rightarrow \begin{pmatrix} (t,x) \\ z \\ p \end{pmatrix} = \begin{pmatrix} s,y+ys+\frac{s^2}{2} \\ y+s \\ \frac{1}{1+s} \begin{pmatrix} 1-y \\ 1 \end{pmatrix} \end{pmatrix} \Rightarrow (s,y) = \left( t,\frac{x-\frac{t^2}{2}}{1+t}\right) \Rightarrow u(t,x) = z(y,s) = y+s = t + \frac{x-\frac{t^2}{2}}{1+t} \]
\end{bsp}
	
\begin{bem} \label{bem_10}
	Für quasilineare PDGL der Form \marginnote{[10]}
	\[ b(u(x),x) \cdot \nabla u(x) + c(u(x),x) = 0 \]
	können die charakteristischen Gleichungen \eqref{eq_2204_3}, \eqref{eq_2204_4} und \eqref{eq_2204_5} vereinfacht werden: $b(z(s),x(s)) \cdot p(s) + c(z(s),x(s)) = 0$ impliziert:
	\begin{align}
		\dot{x}(s) = b(z(s),x(s))	\label{eq_2504_8} \\
		\dot{z}(s) = -c(z(s),x(s))	\label{eq_2504_9} \\
		\dot{p}(s) = -(b_z \cdot p + c_z)\cdot p - b_x p -c_x	\label{eq_2504_10}
	\end{align}
	Die ersten beiden Gleichungen legen bereits $x$ und $z$ fest, sodass \eqref{eq_2504_10} ignoriert werden kann.
\end{bem}
\newpage

%\appendix							% Anhang beginnt hier (Nummerierung mit Buchstaben)
\printindex							% Index ausgeben
\subsection*{Liste der Sätze und Definitionen}
\listtheorems{satz,defn,lemma,anw,korollar}				% listet angegebene Sätze mit Namen und Seitenzahl auf
\end{document}